% ============================================================================
%  front_hms.tex —— 前置部分：封面、译者说明、目录、引言、第一篇开场
%  对应原 tex 第 52–158 行
%  说明：本目录另存有另一项目（arXiv:1106.4789 中译本）的 frontmatter.tex，
%        为免冲突，本文件独立命名为 front_hms.tex。
% ============================================================================

\begin{titlepage}
	\begin{center}

		\LARGE \textbf{论同调镜像对称}\\[0.6cm]

		{\large \textit{About Homological Mirror Symmetry}}\\[1.2cm]

		\Large Alessandro Imparato\footnote{苏黎世联邦理工学院（ETH Zürich）数学系，2022 年夏—秋。本文最后更新于 2023 年 6 月 21 日。}\\[2.4cm]


		\begin{abstract}
			本文讨论 Kontsevich 同调镜像对称猜想的 \textbf{B-侧}。我们首先对导出范畴及其同调代数作一自足的研讨；随后把范围收缩到概形的有界导出范畴，并最终落实到 Calabi--Yau 流形，尤其着重于作为其基础的层论之基本内容及其中的代数特征。最后，我们粗略地讨论同调镜像对称在最低维的若干表现，即椭圆曲线与 $K3$ 曲面情形。

			本文是作者关于同调镜像对称 A-侧的综述 \cite{[Imp21]} 的续篇。
		\end{abstract}

	\end{center}

	\vspace*{4cm}

	\noindent 本文为作者的硕士学位论文，由 Paul Biran 教授指导，提交以取得\textit{苏黎世联邦理工学院应用数学理学硕士（Master of Science ETH in Applied Mathematics）}学位。本文应可供研究生阅读。\\

	\noindent\normalsize\textbf{关键词}：导出范畴、导出函子、三角范畴、拟凝聚层与凝聚层、Calabi--Yau 簇、同调镜像对称、椭圆曲线。

\end{titlepage}


% ----------------------------------------------------------------------------
% 译者说明（非原文内容，为译本所加）
% ----------------------------------------------------------------------------
\newpage
\thispagestyle{plain}
\vspace*{0.5cm}

\begin{center}
\begin{minipage}{0.92\textwidth}
\small
\noindent\textbf{译者说明}

\vspace{0.6em}
\noindent 一、\textbf{底本}。本译本依据 arXiv:2306.13589v2（Alessandro Imparato，\textit{About Homological Mirror Symmetry}，138 页）之作者原始 \LaTeX{} 源码译出，而非据 PDF 图像识别，故全部数学公式与符号均与原稿逐字符一致。

\vspace{0.4em}
\noindent 二、\textbf{学科定位}。原文属纯粹数学（代数几何、同调代数、范畴论），并非量子场论或弦论文献。因此术语一律按数学界约定译为「层」「导出范畴」「凝聚层」等；仅在作者本人援引物理类比（如 D-膜、B-场）时保留物理译名。

\vspace{0.4em}
\noindent 三、\textbf{术语与记号}。关键术语首次出现时以「中文（English）」形式给出；作者本人的记号（如 $\mathsf{D}^{\mathsf{b}}(\mathsf{A})$、$\mathsf{D}^\pi(\widetilde{\mathscr{F}}(X))$）一律沿用，不作改动。若某记号在不同文献中另有含义，译本以 \texttt{[Tutor Note]} 标注。

\vspace{0.4em}
\noindent 四、\textbf{批注}。正文中以 \texttt{\%\,[Tutor Note: ...]} 形式给出的批注为译者所加，非原文内容，涉及：易错推导的符号约定、作者以「显然」跳过之处所省略的证明、以及记号歧义。此类批注在源文件中为 \LaTeX{} 注释，编译后不出现在版面中，仅供阅读源码时参考。

\vspace{0.4em}
\noindent 五、\textbf{图表}。全书有两处带题注的插图（原 PDF 第 134 页、第 137 页）及一处性质蕴含关系示意图（第 79 页），其图像文件均由原 PDF 无损提取，置于 \texttt{figures/} 目录；此外书中的 73 幅交换图本为 \texttt{tikz-cd} 矢量代码，译本保留为可编辑源码，未转为位图。

\vspace{0.4em}
\noindent 六、\textbf{编译}。源文件以 XeLaTeX 引擎编译（中文排版依赖 \texttt{ctex} 宏包），应运行两遍以生成目录与交叉引用；如需索引，须再执行一次 \texttt{makeindex}。
\end{minipage}
\end{center}

\vspace*{1cm}
\begin{flushright}
\small ——译者识
\end{flushright}


% ----------------------------------------------------------------------------
\newpage
\pagenumbering{roman}
\thispagestyle{plain}

\tableofcontents


\newpage

\setcounter{page}{1}

\pagenumbering{arabic}


% ----------------------------------------------------------------------------
\section*{引言}\markboth{引言}{引言}
\addcontentsline{toc}{section}{引言}
\thispagestyle{plain}

\vspace*{1cm}


1994 年，数学家 Maxim Kontsevich 提出了他著名的同调镜像对称猜想（记录于 \cite{[Kon94]}）。这一理论奠基于经典的镜像对称——后者源自众多理论物理学家在研究彼此看似毫无关联的空间之间的对偶性时所作的计算。他们的注意力尤其集中于 \textit{Calabi--Yau 流形}：按超弦理论的假定，这类流形构成我们假想的十维时空的紧致部分；从数学上说，它们是某一类特殊的 Kähler 流形，因而同时具有辛流形与复流形的结构。Kontsevich 的猜想为该领域带来了新的兴趣，它断言：一对镜像 Calabi--Yau 流形 $(X,\overline{X})$ 之间的对偶性，可以用范畴论的抽象语言重述为它们的辛几何「A-侧」与复代数几何「B-侧」之间的关系。

这一纲领性论断随后被 Kenji Fukaya 等研究者逐步精炼，最终形成如下表述：存在三角范畴的等价
\[
\mathsf{D}^\pi(\widetilde{\mathscr{F}}(X))\cong\mathsf{D^b}(\overline{X})\qquad\text{与}\qquad \mathsf{D}^\pi(\widetilde{\mathscr{F}}(\overline{X}))\cong\mathsf{D^b}(X)\,,
\]
其中 $\mathsf{D}^\pi(\widetilde{\mathscr{F}}(X))$ 是 $X$ 的 Fukaya 范畴的\textit{分裂封闭导出范畴}，而 $\mathsf{D^b}(X)$ 是 $X$ 上凝聚层的\textit{有界导出范畴}。同调镜像对称的 A-侧已在 \cite{[Imp21]} 中得到充分讨论；在那里我们论证道，问题归结为理解 $\mathsf{D}^\pi(\widetilde{\mathscr{F}}(X))$（若希望聚焦于镜像伙伴，则相应地是 $\mathsf{D}^\pi(\widetilde{\mathscr{F}}(\overline{X}))$）：（\textit{富化}）\textit{Fukaya 范畴} $\widetilde{\mathscr{F}}(X)$ 由 $X$ 的特殊拉格朗日子流形 $L\subset X$ 构成，每个 $L$ 还配备一个局部系统，而局部系统由一个复向量丛 $E\rightarrow L$ 连同其上的平坦联络 $\nabla$ 组成，$\nabla$ 诱导出酉和乐群，并使曲率成为 B-场的常数倍。我们把这样得到的三元组 $\mathcal{L}=(L,E,\nabla)$ 称为 A\textit{-膜}，以向 D-膜致意——用 D-膜的语汇，上述猜想才能显现出更明显的物理含义——并且我们按照分裂封闭导出范畴的配方，把若干 A-膜串联成有界复形。我们还发现，相应的态射空间可以用 Floer 上同调计算。最重要的是，我们看到 $\mathsf{D}^\pi(\widetilde{\mathscr{F}}(X))$ 只捕捉底层 Calabi--Yau 流形 $X$ 的辛几何。

另一方面，在 \cite[第 7.3 节]{[Imp21]} 中我们只简短地提到了 $\mathsf{D^b}(X)\equiv\mathsf{D^b}(\mathsf{Coh}(X))$ 的定义：与它的 Fukaya 对应物十分相似，我们了解到它同样由几何数据构成，只不过这些数据现在被封装在 $X$ 上的\textit{凝聚层范畴} $\mathsf{Coh}(X)$ 之中，另外还配有一套纯粹同调性的机器，即导出范畴（对于 $A_\infty$-范畴，我们在分裂封闭构造中曾不加分辨地套用了这套机器的类似物，当时甚至不敢自称理解了经典理论！）。倘若我们暂且抛开这些基本素材，试图取巧——即通过强行要求 B-侧长得与 A-侧呈镜像对称，来获得一种同调镜像对偶——那么我们就会被引向如下假定：基本对象应当是 $X$ 的某种全纯嵌入的（偶数维）子流形，并配备全纯向量丛 $E_i$；进而我们会主张，任意一对对象（不妨记作 $(E_i,E_j)$）的相关态射应当属于 $X$ 的、取值于 $\text{Hom}(E_i,E_j)$ 的 Dolbeault 上同调群。这一进路（例如 \cite{[Asp04]} 所讥讽者）当然是幼稚而错误的！理由很简单：按此做法，与 A-侧出现的相比，对象太少，态射也就太少。用物理学的说法，我们缺了 B-膜！

我们此处的工作正是要纠正这一错误直觉。在纠正的过程中我们将发现（如上文已有预告）：支配 B-侧的是一套同样丰富的理论。也就是说，我们需要踏上一段旅程，一方面穿越经典的导出范畴理论（研习它时不必像对付 $A_\infty$-范畴那样顾虑诸多微妙之处），另一方面穿越层与概形的代数几何——后者与 Fukaya 范畴不无相似，各有其自身的困难，也就各有不得不作的妥协（不过说真的，读者多半会发现 B-侧其实是更容易的一侧！）。与上一部一样，我们仍然偏重「大图景」式的把握，而非对每一细枝末节作详尽综述，因为这个主题\textit{极为}广博。\\

据此，本文由两部分构成：第一部分自足地概述导出范畴及其同调代数；第二部分讨论各类空间上的层论。具体说来，第 1 章描述阿贝尔范畴——导出范畴的构件；第 2 章引入导出范畴的语言。在温和的条件下，导出范畴之间还可以得到导出函子，这是第 3 章的内容。第 4 章论证导出范畴是三角范畴的特例，从而恢复了在导出过程中失去的某些良好性质。进入第二部分后，第 5 章回顾层的范畴论基本事实，先讨论拓扑空间上的层，继而讨论代数簇与概形上的层，并顺带温习其代数几何语言。第 6 章把第一部分的工具应用于概形上的凝聚层范畴，研究其上同调，最终把焦点限制到 Calabi--Yau 簇的有界导出范畴。最后，第 7 章通过考察该猜想目前已知的最简单证明——椭圆曲线情形——来结束我们对同调镜像对称的探索，并让 $K3$ 曲面在末尾短暂亮相。 \\

\noindent\textbf{致谢。}我要感谢我的导师 Paul Biran，感谢他抽出时间阅读这篇论文，始终认可我对它极为具体的设想，并在我情绪动荡的一段时期给予我支持。我还要感谢 Will Merry 教授，两年前是他把同调镜像对称这一美妙的主题呈现在我面前。我将永远对您心怀感激。
\vspace*{1cm}

\hfill\textit{谨以此文纪念 William John Merry}



\newpage



\part{同调代数}
\markboth{I\quad 同调代数}{I\quad 同调代数}
\thispagestyle{plain}


\vspace*{2cm}


\noindent 导出范畴位于同调代数这一更广阔的理论之中。同调代数发端于 1940 年代，并凭借 Cartan 与 Eilenberg 1956 年的经典著作 \textit{Homological Algebra} 牢固地立足于数学版图。稍后，在 60 年代初，导出范畴由 Alexander Grothendieck 及其学生 Jean-Louis Verdier 首创，尤其见于后者 1963 年的学位论文（报告于 \cite{[Ver96]}）；正是在那里，导出范畴的公理化催生了三角范畴这一概念。引入导出形式体系的初衷，是为 Grothendieck 在代数几何中的研究提供支撑——这或许也解释了它为何传播得相当缓慢，直到近年才重获推动力。这是一场富有成果的共生：在众多分支学科之中，镜像对称本身的知识也因之推进。

我们将看到，范畴导出过程中的关键想法，在于推广交换代数中环或模可以局部化、从而获得「分式」概念的做法。在那里，只需取「分母」属于任一乘法子集即可，即一个对乘法封闭且含乘法单位元的子集。我们要指出，这一方法在历史上又是受代数簇的局部环之定义所启发（详见第二部分）。类似地，我们将先对\textit{任意}给定范畴，通过形式地可逆化其某一类态射并施加泛性质，来刻画它的导出范畴；只有在此之后，我们才会转向一种更自然的构造进路。

当然，凡有范畴之处，函子必随之而来。导出情形也不例外，人们在那里谈论导出函子。实际上它们率先出现（以其「经典」的上同调面貌），见于 Cartan 与 Eilenberg 的著作，其时导出范畴的概念尚未建立。在那之前，数学家们不得不诉诸谱序列之类更为繁复的工具来研究它们；而对方便程序的需求，在一定程度上正是导出形式体系得以发展的动因！

第一部分的目标，是尽可能完整地介绍理解导出范畴所需的代数机器，并逐步逼近出现在同调镜像对称猜想 B-侧的那个对象 $\mathsf{D^b}(\mathsf{A})$；至于我们偏好的底范畴 $\mathsf{A}$ 究竟是什么，将在第二部分具体交代。确切计划如下：

\begin{itemize}[leftmargin=0.5cm]
	\item 在第 1 章中，我们引入待导出的范畴 $\mathsf{A}$ 必须具备的关键特征，即阿贝尔范畴的那些特征，它们由四条公理整齐地刻画；在任一这样的 $\mathsf{A}$ 上，随后便可定义若干极为方便的概念，例如上同调，它使我们能够真正做同调代数。我们还刻画阿贝尔范畴之间的函子，讨论其正合性。

	\item 在第 2 章中，我们正式启程：先温习上链复形、链映射的定义以及它们的相互作用。然后通过泛性质引入导出范畴，并定性描述其由局部化程序而来的形式结构。遗憾的是，导出范畴\textit{不是}阿贝尔范畴……但我们仍可借助替代工具来分析它们，例如正合三角形与区分三角形，它们恰当地模拟了短正合序列的行为。事实上，我们可以用这些工具把任一导出范畴明确构造为其下层同伦范畴 $\mathsf{H(A)}$ 对拟同构类的局部化。接着我们看到 $\mathsf{A}$ 如何嵌入 $\mathsf{D(A)}$。最后，我们定义 Ext-群并导出其若干有用性质，这些性质将反复出现；同时我们还说明阿贝尔范畴及其单个对象在同调维数方面的含义。

	\item 在第 3 章中，我们把注意力转向联系导出范畴的映射，即右导出函子与左导出函子，它们分别源自左正合与右正合的标准函子。我们说明它们在何种条件下存在——此时它们唯一——以及如何构造它们，并证明它们也是正合的，意即在区分三角形上保持正合。此外我们看到，导出函子在上同调中下降为底层标准范畴之间的经典导出函子，并仍保有长正合序列公理、与复合的相容性等良好性质。我们还考察几个无处不在的函子，并对其导出对应物略作阐明。

	\item 在第 4 章中，我们展开三角范畴的想法（在 \cite[第 3.1 节]{[Imp21]} 中仅顺带提及），回顾 Verdier 的公理。我们进而证明：任一阿贝尔范畴的同伦范畴实际上是三角范畴，并把证明推广到相应的导出范畴；这使关于区分三角形、锥、上同调函子以及诸如此类术语的那种乍看突兀的选择显得合乎情理——它们此后在导出意义与三角意义下都可理解。这也为「导出范畴虽非阿贝尔范畴却仍然如此之好」提供了另一种解释！
\end{itemize}

本部分的参考文献主要是 Gelfand 与 Manin 那本出色的著作 \cite{[GM03]}，偶尔也用到 \cite{[HS97]}、\cite{[KS06]}、\cite{[Mac78]} 与 \cite{[Wei94]}——它们只是数学文献中层出不穷的导出范畴经典处理中的五部。除另有声明外，读者可以放心假定我们取材于 \cite{[GM03]}。



\newpage
\pagestyle{fancy}
